\section{A three-subgroup identity for an \texorpdfstring{$S_5$}{S5} comparison quotient}
\label{sec:sn5-identity}

The previous reflection report establishes the optimal total-multiplicity
factor $n-1$ for $S_n$ in the degrees $3\leq n\leq12$, by positive rational
combinations of characters induced from linear subgroup characters.
Here we isolate a more specific comparison in degree five. Three small
subgroups give an integral decomposition and an explicit product of
Hecke $L$-functions. This is an additional comparison relative to the
reports reviewed \cite{priorreflection}; priority in the full literature has not been established.

\subsection{The exact character identity}

Put $G=S_5$, $t=(12)$, and let $\sigma$ be the nontrivial character of
$\langle t\rangle$. Write $\mathrm{std}$ for the four-dimensional standard
character, so that $\mathrm{std}(g)=\#\mathrm{Fix}(g)-1$.
Use the following subgroups and linear characters:
\[
\begin{array}{c|c|c}
H&\text{generators}&\text{character values on generators}\\ \hline
C=C_6 &(12)(345)&\chi((12)(345))=e^{\pi i/3}\\
V=C_2\times C_2 &(12),(34)&\eta((12))=\eta((34))=-1\\
W=C_5\rtimes C_4 &(12345),(2354)&\xi((12345))=1,\quad\xi((2354))=i
\end{array}
\]
The last group is the affine group of the field with five elements:
its generators $a=(12345)$, $b=(2354)$ satisfy $bab^{-1}=a^2$.
The exact identity is
\begin{equation}
 \operatorname{Ind}_{\langle t\rangle}^{G}\sigma-\mathrm{std}
 =\operatorname{Ind}_{C}^{G}\chi
  +\operatorname{Ind}_{V}^{G}\eta
  +\operatorname{Ind}_{W}^{G}\xi .
\label{eq:v5-identity}
\end{equation}
The three summands have degrees $20,30,6$, respectively.

\paragraph{Proof.}
Order the conjugacy types, suppressing fixed points, as
$1,2,3,22,4,32,5$. Their sizes are
\[
(c_T)_T=(1,10,20,15,30,20,24).
\]
For a linear character $\theta$ of $H$, define the aggregated row
$B_{H,\theta}(T)=\sum_{h\in H\cap T}\theta(h)$.
Direct enumeration gives
\[
\begin{array}{c|rrrrrrr}
 &1&2&3&22&4&32&5\\ \hline
B_{C,\chi}&1&-1&-1&0&0&1&0\\
B_{V,\eta}&1&-2&0&1&0&0&0\\
B_{W,\xi}&1&0&0&-5&0&0&4
\end{array}
\]
For $C$, the two order-three values sum to $-1$, and the two
order-six values sum to $1$. In $W$, its four nonidentity translations
are five-cycles of value $1$; its five elements with multiplier $-1$
are double transpositions of value $-1$; its ten elements of order four
have five values $i$ and five values $-i$, and their sum is zero.
All entries are therefore integers despite the nonreal characters.

The induction formula at type $T$ is
\[
 (\operatorname{Ind}_{H}^{G}\theta)(T)
 =\frac{120}{|H|c_T}\,B_{H,\theta}(T).
\]
The coordinate row paired with the left side of
\eqref{eq:v5-identity}, against an arbitrary real class function, equals
\[
 \frac12\delta_1-\frac12\delta_2
 -\left(\frac{c_T(\#\mathrm{Fix}(T)-1)}{120}\right)_T
 =\left(\frac7{15},-\frac23,-\frac16,0,0,\frac16,\frac15\right).
\]
Entry by entry this is
\[
 \frac16 B_{C,\chi}+\frac14 B_{V,\eta}+\frac1{20}B_{W,\xi}.
\]
The denominators are precisely the subgroup orders, so the induction
formula proves \eqref{eq:v5-identity}. The character values of either
side are $(56,-8,-1,0,0,1,1)$.

\subsection{An entire quotient of four Dedekind zeta functions}

Let $M/K$ be any finite Galois extension of number fields with group
identified with $S_5$. Define fixed fields
\[
 E=M^{\langle t\rangle},\quad P=M^{S_4},\quad
 A=M^C,\quad F=M^V,\quad J=M^W,
\]
where $S_4$ is a point stabilizer. Their degrees over $K$ are
$60,5,20,30,6$, respectively. The characters $\chi,\eta,\xi$ correspond
to finite-order Hecke characters over $A,F,J$ of orders $6,2,4$.
Denote their Hecke $L$-functions by $L_A(s,\chi)$, and so forth.
All $L$-functions here include their canonical local factors, including
ramified primes; the induction identities are exact for these factors.

\paragraph{Theorem.}
The quotient
\begin{equation}
 \mathcal Z(s)=\frac{\zeta_M(s)\zeta_K(s)}{\zeta_E(s)\zeta_P(s)}
\label{eq:sn5-four-zeta}
\end{equation}
extends to an entire function, is nonzero at $s=1$, and has the direct
factorization
\begin{equation}
 \mathcal Z(s)=L_A(s,\chi)L_F(s,\eta)L_J(s,\xi).
\label{eq:sn5-four-zeta-factorization}
\end{equation}
Consequently, at every $\rho\in\mathbb C$,
\begin{equation}
 \operatorname{ord}_{\rho}\zeta_E+
 \operatorname{ord}_{\rho}\zeta_P
 \leq \operatorname{ord}_{\rho}\zeta_M+
       \operatorname{ord}_{\rho}\zeta_K.
\label{eq:sn5-four-zeta-order}
\end{equation}
Orders are signed, with negative order denoting a pole.

\paragraph{Proof.}
Induction from the regular representation of $C_2$ gives
$\operatorname{Ind}_{C_2}^{G}\sigma
 =\operatorname{Reg}_{G}-\operatorname{Ind}_{C_2}^{G}1$.
Also $\mathrm{std}=\operatorname{Ind}_{S_4}^{G}1-1$.
Thus the left side of \eqref{eq:v5-identity} has Artin $L$-function
$\mathcal Z$. Applying Artin induction and multiplicativity gives
\eqref{eq:sn5-four-zeta-factorization}, first for $\Re s>1$ and then
by meromorphic continuation. Every character on the right is nontrivial,
so all three Hecke factors are entire and nonzero at $1$.
The theorem follows, including \eqref{eq:sn5-four-zeta-order}.

\paragraph{Related literature and scope.}
Four-zeta quotients also occur in S.~Gun, S.~Hazra and D.~Sahu,
\href{https://doi.org/10.1007/s00605-025-02060-7}{\emph{On holomorphy and
non-vanishing of Artin $L$-functions}}, Monatshefte f\"ur Mathematik
206 (2025), 853--883 \cite{gun}. Their Theorem~1 concerns solvable Galois groups
and specified normal subgroups. The present $S_5$ calculation concerns
nonnormal subgroups and supplies its own positive character certificate.
This comparison does not establish that the precise quotient above is
new in the full literature.

\subsection{The sign-twisted companion}

Let $\varepsilon$ be the sign character, put
$\tau=\varepsilon\,\mathrm{std}$, and set
$\alpha=(\varepsilon|_C)\chi$. The character $\alpha$ has order three.
On $V$, one has $\varepsilon\eta=1$. On $W$, the generator $a$ is even
and $b$ is odd, so $\varepsilon\xi=\overline\xi$.
Twisting \eqref{eq:v5-identity} therefore gives
\[
 \operatorname{Ind}_{\langle t\rangle}^{G}1-\tau
 =\operatorname{Ind}_{C}^{G}\alpha
   +\operatorname{Ind}_{V}^{G}1
   +\operatorname{Ind}_{W}^{G}\overline\xi.
\]
Hence the meromorphic quotient
\begin{equation}
 \mathcal H(s)=\frac{\zeta_E(s)}{L_K(s,\tau,M/K)}
 =L_A(s,\alpha)\,\zeta_F(s)\,L_J(s,\overline\xi)
\label{eq:sn5-comparison-quotient}
\end{equation}
is holomorphic on $\mathbb C\setminus\{1\}$ and has a simple pole
at $1$. In particular,
\begin{equation}
 \operatorname{ord}_{\rho}L_K(s,\tau,M/K)
 \leq \operatorname{ord}_{\rho}\zeta_E(s)
 \qquad(\rho\ne1).
\label{eq:sn5-order-bound}
\end{equation}
No holomorphy assumption on the individual Artin factor
$L_K(s,\tau,M/K)$ is needed: Brauer's theorem provides meromorphy,
and the product in \eqref{eq:sn5-comparison-quotient} proves the claimed
properties of the quotient. The two nontrivial Hecke factors are
nonzero at $1$, and $\zeta_F$ has a simple pole there.

The classical arithmetic inputs are Artin induction, Brauer meromorphy,
and finite-order Hecke holomorphy and nonvanishing at $1$. Their use in
zero-order comparisons is surveyed by R.~Foote, H.~Ginsberg and
V.~K.~Murty, \href{https://doi.org/10.1090/BULL/1492}{\emph{Heilbronn
characters}}, Bull. Amer. Math. Soc. 52 (2015), 465--496 \cite{heilbronn}.

\subsection{What remains open}

The two quotient theorems give specific zero-order comparisons.
They do not prove that $L_K(s,\tau,M/K)$ is entire:
\eqref{eq:sn5-order-bound} permits a negative left-hand side.
They also leave the zero locations of $\zeta_E$ and $\zeta_M$ open.
The one-sided prime estimate required for a further RH transfer remains
an arithmetic input. Both quotient theorems are independent of RH and
Dirichlet GRH.

The character identity verifies one step in a possible general induction.
For $n\geq4$ put
\[
 U_n=(n-1)\operatorname{Ind}_{C_2}^{S_n}\sigma-\operatorname{Reg}_{S_n}+1,
 \qquad
 V_n=\operatorname{Ind}_{C_2}^{S_n}\sigma-\mathrm{std}_n.
\]
Induction in stages gives the exact recurrence
\[
 U_n=\operatorname{Ind}_{S_{n-1}}^{S_n}U_{n-1}+V_n.
\]
A positive monomial decomposition of every $V_n$ would extend the
previous $U_n$ result to every degree. Equation~\eqref{eq:v5-identity}
supplies the degree-five step; no uniform formula is asserted here.

The accompanying \texttt{check\_sn5\_identity.py} verifies subgroup
orders, every multiplication and character relation in the three
subgroups, the individual cyclotomic character sums, and the untwisted
and sign-twisted identities using exact integer and rational arithmetic.
Its only dependency is Python's standard library.
Inputs are in \texttt{sn5\_certificate.json};
\texttt{sn5\_verification.json} records the successful run.
This supports the written proof and is not a Lean formalization.
